\documentclass[11pt,a4paper]{article}

% ── Packages ─────────────────────────────────────────────────────────────────
\usepackage[utf8]{inputenc}
\usepackage[T1]{fontenc}
\usepackage[french]{babel}
\usepackage{geometry}
\geometry{top=2.5cm,bottom=2.5cm,left=2.5cm,right=2.5cm}
\usepackage{lmodern}
\usepackage{microtype}
\usepackage{graphicx}
\usepackage{booktabs}
\usepackage{tabularx}
\usepackage{array}
\usepackage{xcolor}
\usepackage{fancyhdr}
\usepackage{titlesec}
\usepackage{caption}
\usepackage{float}
\usepackage{amsmath}
\usepackage{amssymb}
\usepackage{hyperref}
\usepackage{enumitem}
\usepackage{tcolorbox}
\usepackage{parskip}

% ── Couleurs ──────────────────────────────────────────────────────────────────
\definecolor{meteoBlue}{HTML}{1D3557}
\definecolor{meteoRed}{HTML}{E63946}
\definecolor{meteoGray}{HTML}{F1FAEE}
\definecolor{meteoCyan}{HTML}{457B9D}

% ── Mise en page ─────────────────────────────────────────────────────────────
\pagestyle{fancy}
\fancyhf{}
\fancyhead[L]{\small\color{meteoBlue}\textbf{Météorage} — Phase 2 : Évaluation des modèles}
\fancyhead[R]{\small\color{meteoBlue}Fiches modèles}
\fancyfoot[C]{\small\color{gray}\thepage}
\renewcommand{\headrulewidth}{0.4pt}

% ── Titres ────────────────────────────────────────────────────────────────────
\titleformat{\section}[block]
  {\Large\bfseries\color{meteoBlue}}
  {\thesection}{1em}{}[\vspace{-0.3em}\rule{\linewidth}{0.4pt}]
\titleformat{\subsection}[block]
  {\large\bfseries\color{meteoCyan}}
  {\thesubsection}{1em}{}
\titleformat{\subsubsection}[block]
  {\normalsize\bfseries\color{meteoBlue}}
  {\thesubsubsection}{1em}{}

\captionsetup{font=small, labelfont=bf, margin=10pt}

\newenvironment{keybox}[1]
  {\begin{tcolorbox}[colback=meteoGray,colframe=meteoBlue,title={\textbf{#1}},
    fonttitle=\small, boxrule=0.6pt, arc=3pt]}
  {\end{tcolorbox}}


% ─────────────────────────────────────────────────────────────────────────────
\begin{document}

% ── Page de titre ─────────────────────────────────────────────────────────────
\begin{titlepage}
  \pagecolor{meteoBlue}
  \color{white}
  \vspace*{2cm}
  \begin{center}
    {\Huge\bfseries Prédiction de fin d'orage\\[0.4em]}
    {\large\bfseries Phase 2 — Fiches d'évaluation des modèles}\\[2em]
    \rule{0.5\linewidth}{0.8pt}\\[1.5em]
    {\large Réponse à l'appel d'offre \textbf{Météorage}}\\[0.5em]
    {\normalsize Comparaison de 5 approches de modélisation\\
    pour l'estimation du temps résiduel d'orage}\\[3em]
    {\normalsize Mars 2026}
  \end{center}
  \vfill
  \begin{center}
    \small\textit{5 modèles $\cdot$ 56\,599 éclairs $\cdot$ 2\,627 sessions d'alerte}
  \end{center}
\end{titlepage}
\nopagecolor

\tableofcontents
\newpage

% ═══════════════════════════════════════════════════════════════════════════════
\section{Vue d'ensemble et comparaison}
% ═══════════════════════════════════════════════════════════════════════════════

\begin{figure}[H]
  \centering
  \includegraphics[width=\linewidth]{figures/comparison_global.pdf}
  \caption{Comparaison globale des 5 modèles : erreur absolue moyenne (gauche) et biais de prédiction (droite).
  Le Neural Hawkes domine nettement avec une MAE de 22 min et un biais quasi nul.}
\end{figure}

\begin{table}[H]
\centering
\renewcommand{\arraystretch}{1.3}
\begin{tabularx}{\linewidth}{l *{5}{>{\centering\arraybackslash}X}}
\toprule
\textbf{Modèle} & \textbf{MAE} & \textbf{RMSE} & \textbf{Médiane} & \textbf{P90} & \textbf{Biais} \\
\midrule
Baseline (30 min)      & 60.26 & 94.05  & 33.32 & 132.48 & $-$52.32 \\
XGBoost AFT            & 57.28 & 90.69  & 33.18 & 149.82 & $-$26.74 \\
Hawkes Classique       & 83.99 & 124.37 & 56.88 & 178.96 & $-$83.82 \\
\textbf{Neural Hawkes (GRU)} & \textbf{22.11} & \textbf{32.50} & \textbf{18.00} & \textbf{40.12} & \textbf{$-$1.61} \\
BNN MC Dropout         & 68.20 & 102.98 & 34.31 & 164.40 & $+$1.42 \\
\bottomrule
\end{tabularx}
\caption{Métriques comparatives sur le jeu de validation (14\,770 éclairs, 526 sessions).
Toutes les valeurs sont en minutes.}
\end{table}

\begin{figure}[H]
  \centering
  \includegraphics[width=\linewidth]{figures/error_distributions.pdf}
  \caption{Distribution des erreurs absolues pour les 4 modèles principaux. Le Neural Hawkes (GRU)
  présente une distribution nettement plus resserrée vers zéro que les trois autres.}
\end{figure}

% ═══════════════════════════════════════════════════════════════════════════════
\newpage
\section{Modèle 1 — Baseline (règle fixe 30 minutes)}
% ═══════════════════════════════════════════════════════════════════════════════

\subsection{Fonctionnement}
La règle métier actuelle consiste à maintenir l'alerte pendant une durée fixe de 30 minutes 
après la détection du dernier éclair cloud-to-ground (CG). Le ``temps restant'' prédit est 
donc constant : $\hat{T} = 30$ min, indépendamment du contexte.

\subsection{Résultats}
\begin{keybox}{Métriques Baseline}
\begin{itemize}[noitemsep]
  \item MAE = 60.26 min — erreur moyenne élevée
  \item Biais = $-$52.32 min — sous-estimation chronique du temps restant
  \item Médiane = 33.32 min — la moitié des prédictions sont erronées de plus de 33 min
\end{itemize}
\end{keybox}

Le biais fortement négatif ($-52$ min) signifie que la règle sous-estime systématiquement 
le temps restant. Pour les sessions longues (plusieurs heures), prédire 30 min est 
catastrophiquement optimiste. À l'inverse, pour les sessions très courtes, la règle 
surestime et maintient l'alerte inutilement.

\subsection{Pistes d'amélioration}
Aucune --- c'est la référence. Son rôle est de quantifier le gain apporté par les modèles.

% ═══════════════════════════════════════════════════════════════════════════════
\newpage
\section{Modèle 2 — XGBoost Survival (AFT)}
% ═══════════════════════════════════════════════════════════════════════════════

\subsection{Fonctionnement}
XGBoost est un algorithme d'ensemble par gradient boosting sur des arbres de décision.
La variante \texttt{survival:aft} (\textit{Accelerated Failure Time}) modélise directement 
$\ln(T)$ comme une variable de régression, où $T$ est le temps avant la fin de l'orage.

\textbf{Features utilisées} (20 au total) :
\begin{itemize}[noitemsep]
  \item Temporelles : temps depuis le début, temps depuis le dernier CG
  \item Fenêtres glissantes (15/30 min) : comptage d'éclairs, ratio IC/CG
  \item Spatiales : distance moyenne, tendance de distance (approche/recul)
  \item Physiques : amplitude, maxis
  \item Cycliques : encodage sin/cos de l'heure et du mois
\end{itemize}

\begin{figure}[H]
  \centering
  \includegraphics[width=0.85\linewidth]{figures/xgboost_importance.pdf}
  \caption{Importance des features (gain) pour XGBoost AFT. Le nombre d'éclairs dans les 15 dernières 
  minutes (\texttt{n\_lightnings\_15m}) domine très largement, indiquant que la densité temporelle 
  récente est le signal le plus informatif.}
\end{figure}

\subsection{Résultats}
\begin{keybox}{Métriques XGBoost AFT}
\begin{itemize}[noitemsep]
  \item MAE = 57.28 min — légère amélioration vs baseline
  \item Biais = $-$26.74 min — réduit de moitié vs baseline
  \item La feature \texttt{n\_lightnings\_15m} domine (gain = 915)
  \item Early stopping à 108 boosting rounds (sur 500 max)
\end{itemize}
\end{keybox}

\subsection{Analyse}
L'amélioration modeste s'explique par le fait que XGBoost traite chaque éclair 
\textbf{indépendamment}. Il ne capture pas la dynamique séquentielle : l'ordre 
et l'espacement des éclairs ne sont exploités qu'indirectement via les features 
de fenêtre glissante.

\subsection{Pistes d'amélioration}
\begin{itemize}[noitemsep]
  \item Ajouter des features séquentielles plus riches (tendances sur 5/10/20 min)
  \item Encoder la géographie (lat/lon relative, azimuth récent)
  \item Tester \texttt{survival:cox} (risques proportionnels) comme alternative
  \item Utiliser un split temporel strict (par date) plutôt que par session
\end{itemize}

% ═══════════════════════════════════════════════════════════════════════════════
\newpage
\section{Modèle 3 — Processus de Hawkes (Classique)}
% ═══════════════════════════════════════════════════════════════════════════════

\subsection{Fonctionnement}
Le processus de Hawkes est un processus ponctuel \textbf{autoexcitant} : 
chaque événement (éclair) augmente momentanément la probabilité d'observer un 
prochain événement. L'intensité conditionnelle s'écrit :
\[
  \lambda(t) = \mu + \sum_{t_i < t} \alpha \cdot \beta \cdot e^{-\beta(t - t_i)}
\]
\begin{itemize}[noitemsep]
  \item $\mu$ = taux de fond (intensité de base)
  \item $\alpha$ = magnitude de l'excitation
  \item $\beta$ = taux de décroissance de l'excitation
\end{itemize}

La fin de l'orage est prédite quand $\lambda(t)$ tombe en dessous d'un seuil fixé.

\begin{figure}[H]
  \centering
  \includegraphics[width=\linewidth]{figures/hawkes_intensity.pdf}
  \caption{Exemple d'intensité conditionnelle du processus de Hawkes. Chaque éclair CG 
  (barres bleues) provoque un pic d'excitation qui décroît exponentiellement. 
  Quand l'intensité passe sous le seuil rouge, l'alerte est levée.}
\end{figure}

\subsection{Résultats}
\begin{keybox}{Paramètres et Métriques Hawkes Classique}
\begin{itemize}[noitemsep]
  \item Paramètres ajustés : $\mu = 0.204$, $\alpha = 0.680$, $\beta = 0.680$
  \item MAE = 83.99 min — le plus mauvais modèle, y compris derrière la baseline
  \item Biais = $-$83.82 min — sous-estime drastiquement le temps restant (coupure prématurée)
  \item Le \textit{branching ratio} $\alpha = 0.68 < 1$ confirme la stationnarité du processus
  (chaque événement engendre en moyenne 0.68 descendant, soit un processus sous-critique)
\end{itemize}
\end{keybox}

\subsection{Analyse}
La performance décevante, marquée par une sous-estimation systématique ($-83$ min), 
s'explique par la \textbf{simplicité du noyau exponentiel}. L'excitation retombe très vite 
sous le seuil de tolérance (ici réglé à $\mu + \varepsilon$), menant le modèle à déclarer la 
fin de l'orage bien trop tôt entre deux rafales lointaines. De plus, le modèle univarié 
ignore l'amplitude et la géométrie de la convection, ne lui laissant aucune indication 
si le front approche ou s'éloigne.

\subsection{Pistes d'amélioration}
\begin{itemize}[noitemsep]
  \item Noyau multivarié : distinguer les éclairs IC et CG
  \item Noyau somme d'exponentielles pour capturer plusieurs échelles temporelles
  \item Marques spatio-temporelles : pondérer l'excitation par la distance et l'amplitude
  \item Seuil adaptatif par aéroport
\end{itemize}

% ═══════════════════════════════════════════════════════════════════════════════
\newpage
\section{Modèle 4 — Neural Hawkes Process (GRU)}
% ═══════════════════════════════════════════════════════════════════════════════

\subsection{Fonctionnement}
Le Neural Hawkes Process remplace le noyau exponentiel paramétrique par un 
\textbf{réseau de neurones récurrent} (GRU). L'idée :
\begin{enumerate}[noitemsep]
  \item Chaque éclair est représenté par un vecteur de features (icloud, amplitude, distance, maxis).
  \item Un GRU accumule l'historique des événements dans un état caché $h_t$.
  \item L'intensité est calculée par : $\lambda(t) = \text{softplus}(W_h \cdot h_{t_{\text{last}}} + w_\Delta \cdot \Delta t + b)$
\end{enumerate}

Contrairement au Hawkes classique, ce modèle :
\begin{itemize}[noitemsep]
  \item Apprend des \textbf{interactions non-linéaires} entre les features
  \item Capture des \textbf{dépendances à long terme} dans la séquence
  \item S'adapte aux \textbf{caractéristiques physiques} de chaque éclair
\end{itemize}

\subsection{Résultats}
\begin{keybox}{Métriques Neural Hawkes (GRU) — \textcolor{meteoRed}{Meilleur modèle}}
\begin{itemize}[noitemsep]
  \item \textbf{MAE = 22.11 min} — 3$\times$ meilleur que la baseline
  \item \textbf{RMSE = 32.50 min} — erreurs extrêmes fortement réduites
  \item \textbf{Médiane = 18.00 min} — la moitié des prédictions à moins de 18 min d'erreur
  \item \textbf{Biais = $-$1.61 min} — prédictions quasi non biaisées
  \item \textbf{P90 = 40.12 min} — 90\% des erreurs sous 40 min
\end{itemize}
\end{keybox}

\subsection{Analyse}
Ce modèle combine le meilleur des deux mondes :
\begin{itemize}
  \item La \textbf{formulation en processus ponctuel} du Hawkes, naturelle pour un flux d'éclairs
  \item La \textbf{puissance d'apprentissage} d'un réseau récurrent (GRU)
\end{itemize}
Il chute nettement les erreurs pour les sessions longues (où les modèles à features fixes échouent)
car il peut ``mémoriser'' la dynamique entière de la séquence.

\subsection{Pistes d'amélioration}
\begin{itemize}[noitemsep]
  \item Remplacer le GRU par un \textbf{Transformer} pour capturer les dépendances globales
  \item Ajouter des features spatiales (azimuth, centroïde mobile)
  \item Entraîner un modèle par aéroport pour capturer les spécificités locales
  \item Intégrer un mécanisme d'attention temporelle (Temporal Attention)
  \item Augmenter la taille du dataset avec des données synthétiques
\end{itemize}

% ═══════════════════════════════════════════════════════════════════════════════
\newpage
\section{Modèle 5 — Bayesian Neural Network (MC Dropout)}
% ═══════════════════════════════════════════════════════════════════════════════

\subsection{Fonctionnement}
Le BNN utilise l'approximation \textbf{Monte Carlo Dropout} : un MLP classique 
(128-64-32 neurones) dont les couches Dropout restent actives pendant l'inférence.
En effectuant $N = 100$ passages forward, on obtient une \textbf{distribution de prédictions} :
\begin{itemize}[noitemsep]
  \item La \textbf{moyenne} des $N$ passages donne la prédiction $\hat{T}$
  \item L'\textbf{écart-type} donne l'incertitude $\sigma$
\end{itemize}

L'intérêt opérationnel est majeur : un contrôleur aérien peut voir non seulement 
``l'orage finira dans 45 min'' mais aussi ``$\pm$ 20 min d'incertitude''.

\begin{figure}[H]
  \centering
  \includegraphics[width=\linewidth]{figures/bnn_uncertainty.pdf}
  \caption{\textbf{Gauche} : incertitude prédite vs erreur réelle. Idéalement, les points 
  devraient suivre la diagonale. \textbf{Droite} : distribution de l'incertitude prédite, 
  centrée autour de 20 minutes.}
\end{figure}

\subsection{Résultats}
\begin{keybox}{Métriques BNN MC Dropout}
\begin{itemize}[noitemsep]
  \item MAE = 68.20 min — performance proche de la baseline en termes d'erreur
  \item Biais = $+$1.42 min — quasi nul (prédictions non biaisées)
  \item Incertitude moyenne = 20.98 min
  \item Calibration : 25.1\% dans 1$\sigma$ (idéal 68.3\%) — sous-calibré
\end{itemize}
\end{keybox}

\subsection{Analyse}
Le BNN avec MC Dropout donne des prédictions \textbf{non biaisées} mais avec une erreur 
élevée. La calibration est insuffisante (25\% dans 1$\sigma$ au lieu de 68\%), 
ce qui signifie que le modèle \textbf{sous-estime son incertitude}. 
Cela vient du fait que le MC Dropout est une approximation grossière de l'inférence 
bayésienne véritable, et que le réseau est trop petit et ne traite pas la séquence.

\subsection{Pistes d'amélioration}
\begin{itemize}[noitemsep]
  \item Utiliser une inférence bayésienne complète (Pyro, TFP) au lieu du MC Dropout
  \item Combiner avec le Neural Hawkes : GRU bayésien pour avoir séquentialité + incertitude
  \item Ajouter une perte de calibration (ex: CRPS, calibration loss)
  \item Augmenter la taille du réseau et le nombre d'epochs
  \item Tester le Concrete Dropout (dropout rate appris)
\end{itemize}

% ═══════════════════════════════════════════════════════════════════════════════
\newpage
\section{Synthèse et recommandation}
% ═══════════════════════════════════════════════════════════════════════════════

\begin{keybox}{Classement final}
\begin{enumerate}
  \item \textbf{Neural Hawkes (GRU)} — Le meilleur sur toutes les métriques (MAE = 22.11, biais $\approx$ 0). 
  Sa formulation en processus ponctuel est naturelle pour les séquences d'éclairs.
  \item \textbf{XGBoost AFT} — Meilleur modèle tabulaire (MAE = 57.28), limité par l'absence 
  de modélisation séquentielle. Facile à déployer en production.
  \item \textbf{Baseline (30 min)} — Référence métier (MAE = 60.26). Dépassée en MAE par 
  XGBoost et Neural Hawkes, mais reste devant le BNN et le Hawkes classique.
  \item \textbf{BNN MC Dropout} — MAE plus élevée que la baseline (68.20 vs 60.26), mais 
  apporte une information complémentaire unique : l'incertitude. Son biais quasi nul (+1.42) 
  et sa capacité à quantifier la confiance en font un candidat intéressant \textit{en complément} 
  d'un modèle plus précis.
  \item \textbf{Hawkes Classique} — Le noyau exponentiel univarié est trop faible (MAE = 83.99), 
  mais le cadre théorique est validé par le succès spectaculaire du Neural Hawkes.
\end{enumerate}
\end{keybox}

\vspace{1em}

La recommandation est de retenir le \textbf{Neural Hawkes Process} comme modèle principal, 
éventuellement enrichi de composantes bayésiennes pour produire des intervalles de confiance 
opérationnels. Un XGBoost AFT pourrait servir de modèle de repli simple pour les 
environnements à faibles ressources de calcul.

\end{document}
